% 题干转写：完成
% 解析转写：完成
% 图形核对：完成
\section{第35届复赛第5题：塞曼效应的经典解释}
\subsection{题干}
钠光 $D$ 线在磁场中会分裂为三条。设原子中的价电子质量为 $m$、电荷量为 $-e$（$e>0$），受到指向原子中心的等效线性回复力
\[
\boldsymbol F_{\mathrm r}=-m\omega_0^2\boldsymbol r,
\]
其中 $\boldsymbol r$ 是电子相对于原子中心的位矢。电子的固有圆频率为 $\omega_0$。现加一沿 $z$ 轴正方向的匀强磁场，磁感应强度参数化为
\[
B=\frac{2m}{e}\omega_L.
\]
求：（1）选择绕磁场方向匀角速转动的参考系，使电子在该系中作简谐运动，导出直角坐标分量方程及其解；（2）将该解变换回实验室参考系；（3）证明原来圆频率为 $\omega_0$ 的谱线分裂为三条，并求弱磁场 $\omega_L\ll\omega_0$ 时相邻谱线的圆频率间隔。已知转动参考系中的科里奥利力为 $\boldsymbol f_{\mathrm c}=-2m\boldsymbol\omega\times\boldsymbol v'$。

\subsection{参考解答整理}
\paragraph{（1）旋转参考系中的简谐运动。} 实验室系中电子满足
\[
m\boldsymbol a=-m\omega_0^2\boldsymbol r-e\boldsymbol v\times\boldsymbol B.
\]
取绕 $z$ 轴转动的参考系 $Ox'y'z'$，并令其角速度为
\[
\boldsymbol\omega=\omega_L\boldsymbol k'.
\]
利用 $\boldsymbol r=\boldsymbol r'$、$\boldsymbol v=\boldsymbol\omega\times\boldsymbol r'+\boldsymbol v'$，离心力和科里奥利力中的速度耦合项相消。电子在旋转系中满足
\[
\ddot x'=-(\omega_0^2+\omega_L^2)x',\qquad
\ddot y'=-(\omega_0^2+\omega_L^2)y',\qquad
\ddot z'=-\omega_0^2z'.
\]
记
\[
\Omega=\sqrt{\omega_0^2+\omega_L^2},
\]
则一般解可写为
\[
x'=A_1\cos(\Omega t+\varphi_1),\qquad
y'=A_2\sin(\Omega t+\varphi_2),\qquad
z'=A_3\cos(\omega_0t+\varphi_3).
\]

\paragraph{（2）变换回实验室系。} 两坐标系之间的旋转变换为
\[
x=x'\cos\omega_Lt-y'\sin\omega_Lt,\qquad
y=x'\sin\omega_Lt+y'\cos\omega_Lt,\qquad z=z'.
\]
将旋转系解代入后，以积化和差公式合并同频项，实验室系中的运动可整理为
\[
\begin{aligned}
x&=A_+\cos(\omega_+t+\varphi_+)+A_-\cos(\omega_-t+\varphi_-),\\
y&=A_+\sin(\omega_+t+\varphi_+)-A_-\sin(\omega_-t+\varphi_-),\\
z&=A_3\cos(\omega_0t+\varphi_3),
\end{aligned}
\]
其中
\[
\boxed{\omega_\pm=\sqrt{\omega_0^2+\omega_L^2}\pm\omega_L.}
\]

\paragraph{（3）谱线分裂与弱场极限。} 上式表明电子的振动包含 $\omega_+$、$\omega_0$、$\omega_-$ 三个圆频率，因此原来的谱线在磁场中分裂为三条。弱磁场条件下，
\[
\sqrt{\omega_0^2+\omega_L^2}
=\omega_0+O\!\left(\frac{\omega_L^2}{\omega_0}\right),
\]
从而
\[
\omega_\pm\simeq\omega_0\pm\omega_L,\qquad
\boxed{\omega_+-\omega_0=\omega_0-\omega_-=\omega_L.}
\]

\subsection{转写依据}
题干、答案公式分别依据 \texttt{CPHO35-Q05.md} 和 \texttt{CPHO35-Q05\_formula\_check\_v1.csv}。原题无题图，故本题无TikZ图件；当前PDF未明示分值，未写入图册。
